\documentclass[10pt,a4paper]{amsart}
\usepackage{amsmath,amssymb,amsthm,graphicx}
% Retained preamble for the Stacks Project chapter {CH} (shared project preamble).
% The chapter file pinned-original-source.tex only carries \input{{preamble}} on its first line,
% so no preamble text lives inside the chapter itself.  The Stacks Project uses one single shared
% preamble for every chapter: the snapshot stages it three times (sources/stacks-chow/preamble.tex,
% sources/stacks-etale-cohomology/preamble.tex and focused-seeds/AGP/sources/stacks-snapshot-20261011/
% preamble.tex) and those copies are identical in size.  The block below is copied mechanically,
% line for line, from sources/stacks-chow/preamble.tex lines 14-18 and 154-217: the Xy-pic setup and
% the native theorem declarations together with the mathematical macros.  Nothing was added, removed
% or redefined; the wrapper class and the amsmath/amssymb/amsthm/graphicx base packages are supplied
% by the assembler.
% For commutative diagrams we use Xy-pic
% Chapter: stacks-cohomology.tex -- Cohomology of Sheaves.
\usepackage[all]{xy}

% We use 2cell for 2-commutative diagrams.
\xyoption{2cell}

% Theorem environments.
%
\theoremstyle{plain}
\newtheorem{theorem}[subsection]{Theorem}
\newtheorem{proposition}[subsection]{Proposition}
\newtheorem{lemma}[subsection]{Lemma}

\theoremstyle{definition}
\newtheorem{definition}[subsection]{Definition}
\newtheorem{example}[subsection]{Example}
\newtheorem{exercise}[subsection]{Exercise}
\newtheorem{situation}[subsection]{Situation}

\theoremstyle{remark}
\newtheorem{remark}[subsection]{Remark}
\newtheorem{remarks}[subsection]{Remarks}

\numberwithin{equation}{subsection}

% Macros
%
\def\lim{\mathop{\mathrm{lim}}\nolimits}
\def\colim{\mathop{\mathrm{colim}}\nolimits}
\def\Spec{\mathop{\mathrm{Spec}}}
\def\Hom{\mathop{\mathrm{Hom}}\nolimits}
\def\Ext{\mathop{\mathrm{Ext}}\nolimits}
\def\SheafHom{\mathop{\mathcal{H}\!\mathit{om}}\nolimits}
\def\SheafExt{\mathop{\mathcal{E}\!\mathit{xt}}\nolimits}
\def\Sch{\mathit{Sch}}
\def\Mor{\mathop{\mathrm{Mor}}\nolimits}
\def\Ob{\mathop{\mathrm{Ob}}\nolimits}
\def\Sh{\mathop{\mathit{Sh}}\nolimits}
\def\NL{\mathop{N\!L}\nolimits}
\def\CH{\mathop{\mathrm{CH}}\nolimits}
\def\proetale{{pro\text{-}\acute{e}tale}}
\def\etale{{\acute{e}tale}}
\def\QCoh{\mathit{QCoh}}
\def\Ker{\mathop{\mathrm{Ker}}}
\def\Im{\mathop{\mathrm{Im}}}
\def\Coker{\mathop{\mathrm{Coker}}}
\def\Coim{\mathop{\mathrm{Coim}}}

% Boxtimes
%
\DeclareMathSymbol{\boxtimes}{\mathbin}{AMSa}{"02}

%
% Macros for moduli stacks/spaces
%
\def\QCohstack{\mathcal{QC}\!\mathit{oh}}
\def\Cohstack{\mathcal{C}\!\mathit{oh}}
\def\Spacesstack{\mathcal{S}\!\mathit{paces}}
\def\Quotfunctor{\mathrm{Quot}}
\def\Hilbfunctor{\mathrm{Hilb}}
\def\Curvesstack{\mathcal{C}\!\mathit{urves}}
\def\Polarizedstack{\mathcal{P}\!\mathit{olarized}}
\def\Complexesstack{\mathcal{C}\!\mathit{omplexes}}
% \Pic is the operator that assigns to X its picard group, usage \Pic(X)
% \Picardstack_{X/B} denotes the Picard stack of X over B
% \Picardfunctor_{X/B} denotes the Picard functor of X over B
\def\Pic{\mathop{\mathrm{Pic}}\nolimits}
\def\Picardstack{\mathcal{P}\!\mathit{ic}}
\def\Picardfunctor{\mathrm{Pic}}
\def\Deformationcategory{\mathcal{D}\!\mathit{ef}}

% --- appended 2026-10-12 (Chief): the project own definitions of the reference/slogan/history
% display environments, copied mechanically from the shared Stacks preamble lines 8-12
% (\usepackage{verbatim} plus the three \newenvironment lines, which the chapter preamble keeps
% verbatim from the project preamble). Nothing added, removed or redefined: these environments
% carry no mathematical content and the project itself typesets them as comments. Without these
% three lines a counted statement whose own text uses them fails with Environment slogan undefined.
\usepackage{verbatim}
\newenvironment{reference}{\comment}{\endcomment}
\newenvironment{slogan}{\comment}{\endcomment}
\newenvironment{history}{\comment}{\endcomment}
\usepackage{cleveref}
\crefname{originalreference}{original reference}{original references}
\Crefname{originalreference}{Original reference}{Original references}
\makeatletter
\newlabel{derived-remark-map-into-derived-limit-truncations}{{0H72}{source}{}{sourcepointer-1}{}}
\newlabel{derived-remark-map-into-derived-limit-truncations@cref}{{[originalreference][0][]0H72}{[1][0][]source}{}{}{}}
\makeatother
\crefname{originallemma}{original lemma}{original lemmas}
\Crefname{originallemma}{Original lemma}{Original lemmas}
\makeatletter
\newlabel{lemma-cohomology-derived-limit-injective}{{0D61}{source}{}{sourcepointer-2}{}}
\newlabel{lemma-cohomology-derived-limit-injective@cref}{{[originallemma][0][]0D61}{[1][0][]source}{}{}{}}
\makeatother
\crefname{originalreference}{original reference}{original references}
\Crefname{originalreference}{Original reference}{Original references}
\makeatletter
\newlabel{derived-remark-truncation-distinguished-triangle}{{08J5}{source}{}{sourcepointer-3}{}}
\newlabel{derived-remark-truncation-distinguished-triangle@cref}{{[originalreference][0][]08J5}{[1][0][]source}{}{}{}}
\makeatother
\title{E -> R lim tau\_\{>=-n\}E is an isomorphism under per-point vanishing}
\author{The Stacks Project Authors}
\date{stacks-project:cohomology.tex:local-snapshot-20261011; GNU FDL 1.2 or later}
\begin{document}
\maketitle
\noindent\textit{Editorial note.} All delivered mathematical text is the author's own, copied verbatim from the pinned Stacks Project chapter file cohomology.tex; no statement or proof step was rewritten, reordered or removed.\par\medskip
\noindent\textit{Editorial note.} The counted span cites the lemmas recorded in assembly.external\_references with canonical tags as navigation pointers.\par\medskip
\noindent\textit{Allowed standard prerequisites (editorial).}\begin{itemize}
\item Truncation triangles. Derived Categories, Remark derived-remark-truncation-distinguished-triangle. Distinguished triangles H\textasciicircum{}\{-n-1\}(E)[n+1] -> K\_\{n+1\} -> K\_n -> H\textasciicircum{}\{-n-1\}(E)[n+2].
\item Stalkwise injectivity criterion. Cohomology of Sheaves, Lemma cohomology-lemma-cohomology-derived-limit-injective. Injectivity of H\textasciicircum{}m(R lim K\_n)\_x -> H\textasciicircum{}m(K\_\{n(x)\})\_x from the two per-point assumptions.
\end{itemize}\medskip
\ifdefined\hypertarget\hypertarget{sourcepointer-1}{}\fi
\noindent\textit{Source reference.} Original reference 0H72: sources/stacks-chow/tags.txt, source lines 21442--21442. Cross-chapter standard prerequisite cited by the author inside the counted span; it is used as a pointer and its own proof is not part of this item.\par\medskip
\ifdefined\hypertarget\hypertarget{sourcepointer-2}{}\fi
\noindent\textit{Source reference.} Original lemma 0D61: sources/stacks-chow/tags.txt, source lines 16506--16506. Cross-chapter standard prerequisite cited by the author inside the counted span; it is used as a pointer and its own proof is not part of this item.\par\medskip
\ifdefined\hypertarget\hypertarget{sourcepointer-3}{}\fi
\noindent\textit{Source reference.} Original reference 08J5: sources/stacks-chow/tags.txt, source lines 10838--10838. Cross-chapter standard prerequisite cited by the author inside the counted span; it is used as a pointer and its own proof is not part of this item.\par\medskip
\begin{lemma}
\label{lemma-is-limit-per-point}
Let $(X, \mathcal{O}_X)$ be a ringed space. Let $E \in D(\mathcal{O}_X)$.
Assume that for every $x \in X$ there exist
a function $p(x, -) : \mathbf{Z} \to \mathbf{Z}$ and
a fundamental system $\mathfrak{U}_x$ of open neighbourhoods of $x$
such that
$$
H^p(U, H^{m - p}(E)) = 0 \text{ for }
U \in \mathfrak{U}_x \text{ and } p > p(x, m)
$$
Then the map $E \to R\lim \tau_{\geq -n} E$ of
Derived Categories, Remark
\ref{derived-remark-map-into-derived-limit-truncations}
is an isomorphism in $D(\mathcal{O}_X)$.
\end{lemma}

\begin{proof}
Set $K_n = \tau_{\geq -n}E$ and $K = R\lim K_n$.
The canonical map $E \to K$
comes from the canonical maps $E \to K_n = \tau_{\geq -n}E$.
We have to show that $E \to K$ induces an isomorphism
$H^m(E) \to H^m(K)$ of cohomology sheaves. In the rest of the
proof we fix $m$. If $n \geq -m$, then
the map $E \to \tau_{\geq -n}E = K_n$ induces an isomorphism
$H^m(E) \to H^m(K_n)$.
To finish the proof it suffices to show that for every $x \in X$
there exists an integer $n(x) \geq -m$ such that the map
$H^m(K)_x \to H^m(K_{n(x)})_x$ is injective. Namely, then
the composition
$$
H^m(E)_x \to H^m(K)_x \to H^m(K_{n(x)})_x
$$
is a bijection and the second arrow is injective, hence the
first arrow is bijective. Set
$$
n(x) = 1 + \max\{-m, p(x, m - 1) - m, -1 + p(x, m) - m, -2 + p(x, m + 1) - m\}.
$$
so that in any case $n(x) \geq -m$. Claim: the maps
$$
H^{m - 1}(U, K_{n + 1}) \to H^{m - 1}(U, K_n)
\quad\text{and}\quad
H^m(U, K_{n + 1}) \to H^m(U, K_n)
$$
are isomorphisms for $n \geq n(x)$ and $U \in \mathfrak{U}_x$.
The claim implies conditions
(1) and (2) of Lemma \ref{lemma-cohomology-derived-limit-injective}
are satisfied and hence implies the desired injectivity.
Recall (Derived Categories, Remark
\ref{derived-remark-truncation-distinguished-triangle})
that we have distinguished triangles
$$
H^{-n - 1}(E)[n + 1] \to
K_{n + 1} \to K_n \to H^{-n - 1}(E)[n + 2]
$$
Looking at the associated long exact cohomology sequence the claim follows if
$$
H^{m + n}(U, H^{-n - 1}(E)),\quad
H^{m + n + 1}(U, H^{-n - 1}(E)),\quad
H^{m + n + 2}(U, H^{-n - 1}(E))
$$
are zero for $n \geq n(x)$ and $U \in \mathfrak{U}_x$.
This follows from our choice of $n(x)$
and the assumption in the lemma.
\end{proof}

\end{document}
